\section*{Chapter \ref{ch:g11:seismic-structure} --- Seismic Waves Reveal the Interior}

\begin{solution}{exo:g11:seismic-structure:1}
A \omterm{def:g11:seismic-structure:body-waves}{P wave} compresses and stretches the rock, which moves along the direction
of travel; an \omterm{def:g11:seismic-structure:body-waves}{S wave} shears it, and the rock moves at right angles to that
direction. \omterm{def:g11:seismic-structure:body-waves}{P waves} cross solids and liquids; \omterm{def:g11:seismic-structure:body-waves}{S waves} cross only solids. The \omterm{def:g11:seismic-structure:body-waves}{P
wave} is recorded first, because it is the faster of the two.
\end{solution}

\begin{solution}{exo:g11:seismic-structure:2}
$40 \times 111.2 = \qty{4448}{km}$, about \qty{4450}{km}. $8900/111.2
\approx \ang{80}$.
\end{solution}

\begin{solution}{exo:g11:seismic-structure:3}
\qty{456}{s} is 7~minutes 36~seconds; \qty{822}{s} is 13~minutes
42~seconds. The S--P delay is $822 - 456 = \qty{366}{s}$, that is 6~minutes
6~seconds.
\end{solution}

\begin{solution}{exo:g11:seismic-structure:4}
$\sin i_c = 6.0/8.0 = 0.75$, so $i_c \approx \ang{48.6}$.
\end{solution}

\begin{solution}{exo:g11:seismic-structure:5}
The outer core is $5150 - 2891 = \qty{2259}{km}$ thick. The core's radius is
$6371 - 2891 = \qty{3480}{km}$, and $3480/6371 \approx 0.55$: a little more
than half the Earth's radius.
\end{solution}

\begin{solution}{exo:g11:seismic-structure:6}
6~minutes 6~seconds is \qty{366}{s}: the earthquake is \ang{40} away, that is
$40 \times 111.2 \approx \qty{4450}{km}$. A single delay gives only the
distance: the epicentre can be anywhere on a circle of that radius around the
station. Two more stations are needed to find the one point at the right
distance from all three.
\end{solution}

\begin{solution}{exo:g11:seismic-structure:7}
$i_c \approx \ang{48.6}$ and $\cos i_c \approx 0.661$. Direct wave: $200/6.0
\approx \qty{33.3}{s}$. \omterm{ex:g11:seismic-structure:sp}{Head wave}: $200/8.0 + 2 \times 35 \times 0.661/6.0
\approx 25.0 + 7.7 = \qty{32.7}{s}$. The \omterm{ex:g11:seismic-structure:sp}{head wave} arrives first, about
\qty{0.6}{s} earlier.
\end{solution}

\begin{solution}{exo:g11:seismic-structure:8}
$\tan i_c \approx 1.134$: the critical distance is $2 \times 35 \times 1.134
\approx \qty{79}{km}$. The \omterm{prop:g11:seismic-structure:crossover}{crossover distance} is $2 \times 35 \times
\sqrt{14.0/2.0} = 70\sqrt 7 \approx \qty{185}{km}$. The station of
\cref{exo:g11:seismic-structure:7} lies at \qty{200}{km}, beyond the
crossover: the \omterm{ex:g11:seismic-structure:sp}{head wave} arrives first, as found.
\end{solution}

\begin{solution}{exo:g11:seismic-structure:9}
$\sin i_c = 6.5/8.0 = 0.8125$, $i_c \approx \ang{54.3}$ and $\cos i_c \approx
0.583$. \omterm{prop:g11:seismic-structure:head-wave}{Intercept time}: $2 \times 6 \times 0.583/6.5 \approx \qty{1.1}{s}$.
\omterm{prop:g11:seismic-structure:crossover}{Crossover distance}: $12\sqrt{14.5/1.5} \approx \qty{37}{km}$. Under a
continent the crossover lies near \qty{165}{km}, so stations a few tens of
kilometres apart catch both lines; under the ocean the \omterm{ex:g11:seismic-structure:sp}{head wave} already
arrives first beyond about \qty{37}{km}, and the direct-wave line is only a
few tens of kilometres long: the stations must be a few kilometres apart to
draw it.
\end{solution}

\begin{solution}{exo:g11:seismic-structure:10}
$\sin r = (8.06/13.7)\sin\ang{30} \approx 0.294$, so $r \approx \ang{17.1}$:
the ray is bent towards the perpendicular, that is towards the centre of the
Earth. It therefore crosses the core along a long path close to its middle,
and comes back out far round the other side: P rays that enter the core
emerge much farther from the earthquake than the rays that just miss it.
\end{solution}

\begin{solution}{exo:g11:seismic-structure:11}
$12\,742/13.7 \approx \qty{930}{s}$, about 15.5~minutes, whereas the wave
takes \qty{1210}{s}, about 20~minutes: it is some 280~seconds late. Even at
the highest speed of the mantle the wave would arrive much earlier, so the
middle of the Earth must slow it down. Oldham concluded that the Earth has a
central core in which \omterm{def:g11:seismic-structure:body-waves}{P waves} travel more slowly than in the mantle.
\end{solution}

\begin{solution}{exo:g11:seismic-structure:12}
(1) By the proposition, $t_i = 2h\cos i_c/v_1$, and $\cos i_c = \sqrt{1 -
v_1^2/v_2^2} = \sqrt{v_2^2 - v_1^2}/v_2$. So
\[
  t_i = \frac{2h\sqrt{v_2^2 - v_1^2}}{v_1 v_2},
  \qquad\text{hence}\qquad
  h = \frac{t_i\,v_1 v_2}{2\sqrt{v_2^2 - v_1^2}}.
\]
(2) $h = 6.23 \times 6.2 \times 8.1/\bigl(2\sqrt{65.61 -
38.44}\bigr) \approx 312.9/10.43 \approx \qty{30.0}{km}$.
\end{solution}

\begin{solution}{exo:g11:seismic-structure:13}
(1) $(3480/6371)^3 \approx 0.163$: the core is about \qty{16}{\percent} of
the volume. (2) The mean density is the volume-weighted mean: $5.52 = 0.163
\rho_c + 0.837 \times 4.45$, so $\rho_c \approx (5.52 - 3.72)/0.163 \approx
\qty{11.0}{g/cm^3}$. (3) This is more than \qty{9.9}{g/cm^3}, as it should
be, since the density of the core increases towards the centre (to
\qty{13.1}{g/cm^3}); it is about four times the density of the crust,
\qty{2.6}{g/cm^3}. The core is made of much denser material than the rocks
of the surface.
\end{solution}

\begin{solution}{exo:g11:seismic-structure:14}
(1) $105 \times 111.2 \approx \qty{11700}{km}$ and $143 \times 111.2 \approx
\qty{15900}{km}$. (2) $(\cos\ang{105} - \cos\ang{143})/2 \approx (-0.259 +
0.799)/2 \approx 0.27$: about \qty{27}{\percent} of the surface, a ring around
the point opposite the earthquake. (3) The \omterm{def:g11:seismic-structure:shadow-zone}{shadow zone} is a range of
distances from the epicentre: it moves with the earthquake. Each earthquake
has its own ring.
\end{solution}

\begin{solution}{exo:g11:seismic-structure:15}
(1) $h = 9.5 \times 6.3 \times 8.1/\bigl(2\sqrt{65.61 - 39.69}\bigr) \approx
484.8/10.18 \approx \qty{48}{km}$. (2) $\sin i_c = 6.3/8.1 \approx 0.778$,
$\cos i_c \approx 0.629$, and $t_i = 2 \times 30 \times 0.629/6.3 \approx
\qty{6.0}{s}$. (3) The crust under the range is about \qty{48}{km} thick,
thicker than the continental average of about \qty{41}{km} and close to the
average of mountain belts: high ranges stand on thick crust.
\end{solution}

\begin{solution}{pb:g11:seismic-structure:1}
\textbf{1.} The Pg points lie on a straight line; the Pn points lie on
another straight line, less steep.

\textbf{2.} The earthquake is at a very shallow depth beneath the epicentre:
a station at distance zero would record the wave at once.

\textbf{3.} $(32.3 - 3.3)/(200 - 20) = 29.0/180 \approx \num{0.161}$~\unit{s}
per kilometre.

\textbf{4.} $v_1 = 180/29.0 \approx \qty{6.2}{km/s}$.

\textbf{5.} $3.3 + 280 \times 0.161 \approx \qty{48.4}{s}$.

\textbf{6.} $(43.2 - 18.6)/(300 - 100) = 24.6/200 = \num{0.123}$~\unit{s}
per kilometre, so $v_2 = 200/24.6 \approx \qty{8.13}{km/s}$.

\textbf{7.} It is close to \qty{8.0}{km/s}: beneath the crust lies the rock
of the uppermost mantle.

\textbf{8.} $\sin i_c = 6.21/8.13 \approx 0.763$, so $i_c \approx
\ang{49.8}$.

\textbf{9.} The \omterm{ex:g11:seismic-structure:sp}{head wave} goes down at the angle $i_c$ and comes back up at
the same angle. The first point where it can return to the surface lies
$h\tan i_c$ beyond the point where it reached the \omterm{def:g11:seismic-structure:moho}{Moho}, itself $h\tan i_c$
from the source: no station closer than $2h\tan i_c$ can receive it.

\textbf{10.} A seismometer records every wave that reaches it, not only the
first; the Pn wave arrives after the Pg wave but is still seen as a separate
wave on the record.

\textbf{11.} At \qty{160}{km}, Pg (\qty{25.8}{s}) comes before Pn
(\qty{26.0}{s}); at \qty{180}{km}, Pn (\qty{28.5}{s}) comes before Pg
(\qty{29.0}{s}). The change happens between these two stations.

\textbf{12.} Pg: $t = 3.3 + 0.1611\,(x - 20) \approx 0.08 + 0.1611\,x$. Pn:
$t = 18.6 + 0.123\,(x - 100) = 6.30 + 0.123\,x$ ($x$ in kilometres, $t$ in
seconds).

\textbf{13.} $0.08 + 0.1611\,x = 6.30 + 0.123\,x$ gives $x \approx
6.22/0.0381 \approx \qty{163}{km}$.

\textbf{14.} Pn, after $6.30 + 0.123 \times 250 \approx \qty{37.1}{s}$.

\textbf{15.} $t_i \approx \qty{6.3}{s}$.

\textbf{16.} $\cos i_c \approx 0.646$, so $h = 6.21 \times 6.3/(2 \times
0.646) \approx \qty{30.3}{km}$.

\textbf{17.} $x_c = 2 \times 30.3 \times \sqrt{14.34/1.92} \approx 60.6 \times
2.73 \approx \qty{165}{km}$, within \qty{2}{km} of the \qty{163}{km} of
question~13.

\textbf{18.} $\tan i_c \approx 1.18$: $2 \times 30.3 \times 1.18 \approx
\qty{72}{km}$. The nearest station that records Pn is at \qty{80}{km}, and
the one at \qty{60}{km} does not: the two agree.

\textbf{19.} The times are read to \qty{0.1}{s}, and an error of
\qty{0.1}{s} on $t_i$ alone changes $h$ by about \qty{0.5}{km}; the speeds,
taken from slopes, are uncertain too, and the earthquake was not exactly at
the surface. The crust is also not perfectly uniform. A precision of a
kilometre or two is the most the data allow.

\textbf{20.} \textbf{The \omterm{def:g11:seismic-structure:moho}{Moho} lies about \qty{30}{km} beneath this region}:
the crust here is on the thin side for a continent, whose crust averages about
\qty{41}{km}; \qty{30}{km} is typical of regions where the crust has been
stretched.
\end{solution}
